%% sections/introduction.tex — Step B11
\section{Introduction}
\label{sec:intro}

Lung cancer remains the leading cause of cancer-related mortality worldwide,
with non-small cell lung cancer (NSCLC) accounting for approximately 85\% of
cases \cite{Siegel2024}.
The introduction of immune checkpoint inhibitors (ICIs) targeting the
PD-1/PD-L1 axis has transformed treatment for advanced NSCLC, producing
durable responses in a subset of patients \cite{Reck2016,Borghaei2015}.
However, only 20--30\% of unselected patients achieve an objective response,
and a substantial proportion experience primary or acquired resistance
\cite{Rizvi2015}.
Reliable pre-treatment identification of likely responders remains an
unresolved clinical need, as it would allow clinicians to prioritise ICI
therapy for patients most likely to benefit and to consider alternative or
combination strategies for those unlikely to respond.

PD-L1 tumour proportion score (TPS) and tumour mutational burden (TMB) are
the most widely used biomarkers for ICI response prediction, but both have
well-documented limitations: PD-L1 expression is spatially heterogeneous and
assay-dependent, while TMB alone provides only modest discriminative
performance \cite{Herbst2016,Rizvi2015}.
This has motivated growing interest in multimodal approaches that integrate
radiomics from CT imaging, pathomics from digitised histology, genomic
alterations, and clinical variables into a single predictive model
\cite{Bodalal2019,Dercle2020,Trebeschi2019,Cheerla2019}.
Such approaches are appealing because they can capture complementary aspects
of tumour biology and host physiology that no single modality fully
represents; however, they introduce new methodological challenges, including
how to weight modalities of differing reliability and how to handle missing
modalities without discarding patients \cite{Ma2021}.
The Dynamic Attention Multimodal (DyAM) framework
\cite{Vanguri2022} addresses these challenges through a learned, per-patient
attention mechanism that adaptively weights each available modality and
naturally accommodates missing data via masking, but its original
formulation used only a cooperative (softmax/L1-normalised) attention
mechanism and represented clinical variables as raw numeric features.

A largely unexplored question in this setting is how structured clinical
variables --- demographics, laboratory values, performance status --- should
be represented for fusion with imaging- and molecular-derived features.
The conventional approach treats each clinical variable as an independent
numeric input, requiring the model to learn inter-variable relationships
(e.g., the joint prognostic implication of advanced age, low albumin, and
poor performance status) from a relatively small training set.
Recent advances in natural language processing (NLP) have shown that
pretrained sentence encoders can represent complex, multi-attribute entities
as dense embeddings that preserve semantic similarity
\cite{Reimers2019,Wang2020MiniLM}, and domain-specific clinical language
models have been developed for free-text clinical notes
\cite{Alsentzer2019}.
However, whether converting \emph{structured, tabular} clinical variables
into natural-language sentences and encoding them with a general-purpose
sentence transformer can improve multimodal fusion performance --- as
opposed to encoding free-text clinical narratives --- has not, to our
knowledge, been systematically evaluated for ICI response prediction.

A second open question concerns the form of the attention mechanism used to
combine modalities.
Cooperative attention mechanisms, which normalise modality weights to sum to
one via softmax or L1 normalisation, implicitly assume that all modalities
contribute complementary, additive information.
An alternative, competitive formulation --- in which each modality's weight
is determined relative to the others via a one-vs-others (OvO) comparison
--- may better capture settings where modalities are partially redundant and
the most informative modality should dominate.
Whether such a competitive mechanism improves performance, and under what
modality configurations, remains untested in the multimodal ICI response
prediction literature.

In this study, we extend the DyAM framework with two contributions and
evaluate them on a discovery cohort of 247 patients with advanced NSCLC
treated with ICIs.
First, we introduce an OvO competitive attention variant
(\texttt{AttentionMatrixOvO}) and compare it against the original cooperative
attention across 20 modality combinations using 10-fold cross-validation,
with a confirmatory repeated-seed evaluation (5~seeds~$\times$~10-fold) across
21 modality combinations to test the reproducibility of any observed
difference.
Second, we convert 13 numeric clinical variables into natural-language
sentences and encode them with a pretrained sentence transformer
(MiniLM-L6-v2), comparing this NLP-based representation against raw numeric
encoding, a domain-specific clinical language model (BioClinBERT), and a
PCA-compressed variant, evaluated using both binary classification (AUC) and
survival analysis (Cox regression, time-dependent AUC, Harrell's C-index, and
Kaplan--Meier stratification).
Together, these contributions aim to clarify (i) whether competitive
attention offers a reproducible accuracy advantage over cooperative
attention in multimodal fusion, or whether the two are functionally
interchangeable at this sample size, and (ii) whether NLP-based encoding of
structured clinical data provides a practical, fine-tuning-free improvement
over conventional numeric representations for both response prediction and
survival risk stratification.
